\documentclass{article}
\usepackage{pathfinder-readable}

\title{When Learning Has a Price: Participation in a Resource Allocation Mechanism}

\begin{document}
\maketitle

\begin{abstract}
Q studies a resource allocation mechanism whose equilibrium balances transfers and gives each agent
a nonnegative payoff; P proposes charging language-model auction traders for the tokens used in
their decisions. We asked what would happen to Q's participation claim if its learning samples also
carried such a charge. Under that extra assumption, the peers found an accounting limit, two
participation examples, and a finite limit on the length of Q's stated sampling schedule. P does not
test Q's mechanism, and Q does not identify its samples with token-generating calls. The thread
paused because the conditional bounds do not yet give a new mechanism or a finite-budget accuracy
result.
\end{abstract}

\section{Introduction}

Q, by Garjani, Shokri and Kebriaei, allocates a shared resource among agents with private, uncertain
valuations \cite{Q}. Each agent values its own allocation, and the valuations are strongly concave.
Q designs quadratic transfers so that the game's equilibrium gives the welfare-maximising
allocation, balances the transfers, and leaves each agent at least as well off as taking no
allocation. It also gives a decentralised learning algorithm and proves mean-square convergence
under its stated conditions.

P is a local manuscript about a different setting: buyers and sellers represented by language models
in a fixed-value double auction \cite{P}. It proposes charging for generated tokens when traders
decide whether and how to quote. The proposal separates gains from completed trades from the cost of
the calls used to reach them; no combined charged-auction experiment has been run. P calls two
papers in its own source discussion ``Q'' and ``P'', but those labels refer to other works, not to
the assigned Q and P here.

The possible connexion was the cost of reaching an outcome. Q's participation result concerns the
payoff at an ideal equilibrium after learning. P makes the cost of decisions explicit, including a
charge for the first scheduled decision even if a trader withdraws. If Q's samples required paid,
token-generating calls, an agent would face that cost before receiving its eventual allocation
benefit. This is a proposed extension of Q, not a property established by either paper.

\section{What was done}

The peers first checked which papers were actually paired. They then compared Q's equilibrium payoff
with P's cost accounting. They derived a condition on total value and total cost, before testing
whether it could fail within Q's permitted valuations. One test used agents with no positive gain
from an allocation. A second used Q's explicit transfer rule to show that positive total value need
not pay every participant's bill.

Next, the peers examined the number of samples in Q's convergence theorem. They asked how long its
schedule could run if every sample had a positive cost and only one allocation payoff followed
learning. They also checked whether P's buyer--seller matching met Q's assumptions, and whether a
fixed charge would alter Q's stage-game choices. These checks narrowed the question to entry and
learning choices made before a charge is incurred.

\section{Results}

Let \(g_n\) be agent \(n\)'s gross gain, \(c_n\) its metered cost, and \(t_n\) its transfer, counted
as a payment when positive. Let \(\lambda>0\) convert cost units to gain units. Write \(G=\sum_n
g_n\) for total gross gain and \(C=\sum_n c_n\) for total cost. If transfers balance, so that
\(\sum_n t_n=0\), the agents' payoffs \(u_n=g_n-\lambda c_n-t_n\) satisfy
\[
\sum_n u_n=G-\lambda C.
\]
With zero outside payoffs, every agent can be willing to participate only if \(G\geq\lambda C\).
This is an accounting condition. It says nothing by itself about incentives or whether agents can
learn the needed transfers. For P's stated symmetric auction schedule, gross surplus cannot exceed
\(7.5\) in a round. Thus \(\lambda C>7.5\) rules out both balanced transfers and nonnegative payoffs
for every trader in that round \cite{P}.

The first example applies to Q's full valuation class. Give every agent an allocation choice in
\([0,1]\) and valuation \(V_n(x)=-\alpha x^2/2\), where \(\alpha>0\), \(x\) is that agent's
allocation, and \(n\) labels the agent. These valuations meet Q's concavity and zero-at-zero
conditions. The best total gross value is zero. If entry requires even one positive-cost sample
before exit is possible, total net value is negative, so balanced transfers cannot give every
entrant a nonnegative payoff. Emmy's derivation and Ada's independent check establish this
conditional obstruction. A free choice before the first call, outside funding, or a narrower set of
valuations could avoid this example; none alone proves a useful welfare outcome.

The second example shows why a positive net total is insufficient for Q's particular transfer. With
two agents, allocations in \([0,1]\), and capacity \(2\), take \(V_1(x)=-x^2/2\) and
\(V_2(x)=x-x^2/2\). Q's rule admits an equilibrium with allocations \((0,1)\), zero price messages
and zero transfers. Agent 1 gets no gross benefit. If that agent must pay \(0.1\) for learning, its
payoff becomes \(-0.1\), although total net value is \(0.5-0.1=0.4\). Ada checked the equilibrium
against Q's explicit payment rule: at zero rival price, a positive own price only adds a nonnegative
charge, and each agent's stated allocation maximises its valuation. The example challenges a
cost-aware extension of Q's participation claim, not Q's original claim.

Q's Theorem 4 prescribes \(Q_i=\lceil C_b/\eta^{2(i+1)}\rceil\) samples per agent at iteration
\(i\), where \(C_b>0\) is a theorem constant and \(0<\eta<1\). Suppose there are \(N\) agents, every
sample costs at least \(d>0\) in gain units, and the agents receive one allocation payoff after
\(I\) iterations.
The total cost is then at least
\[
dN\sum_{i=0}^{I-1}Q_i
\geq \frac{dNC_b(\eta^{-2I}-1)}{1-\eta^2}.
\]
Let \(M\) be the largest total gross valuation over Q's compact allocation sets. It is finite.
Balanced transfers and nonnegative payoffs for all entrants therefore require
\[
I\leq
\frac{\log\!\left(1+M(1-\eta^2)/(dNC_b)\right)}
{2\log(1/\eta)}.
\]
Ada and Emmy derived this limit independently. It is a necessary cap on iterations under the stated
cost and single-payoff assumptions. It gives no bound on the allocation error at that stopping
point.

\section{Objections and limits}

P's auction cannot simply inherit Q's transfer theorem. Q adds valuations that each depend on one
agent's continuous allocation. In a simple trade needing both a buyer and seller, a gain \(h>0\)
appears only when both participate. Its dependence on their two choices cannot be written as a sum
of separate buyer and seller values. Binary participation also lies outside Q's convex allocation
sets. Emmy checked this mismatch using the positive cross-difference of the joint trade gain.

A sampling bill fixed independently of an agent's messages leaves its best response in Q's stage
game unchanged once the bill has been paid. The consequential choices occur earlier: whether to
enter, buy samples, continue, or leave. P's proposed withdrawal happens on a charged first call, so
it is too late to give the costless pre-call choice needed in the zero-gain example. Q's simulation
text also states a sample count \(\lceil0.96^{i+1}\rceil=1\) at every iteration, unlike the growing
count in its convergence theorem. Its plots therefore do not settle the paid-sampling question.

The peers checked adjacent work on costly reporting \cite{osborne}, data acquisition \cite{chen},
computation on costly data \cite{delegating}, and budget-constrained mechanism design
\cite{budgetagents}. These sources narrow any novelty claim; the checks did not establish novelty
for this particular pairing. The ledger records no experiment in which Q's samples were charged, and
no finite-budget mechanism with an accuracy or welfare guarantee was produced.

\section{Why the thread stopped}

The verifier accepted the accounting bound, examples and iteration cap under their stated
assumptions. It paused the thread because the central bridge remains unproved: Q's samples have not
been shown to incur P-style token costs. The results also stop short of a new incentive or accuracy
theorem. The smallest useful restart is to specify a paid-sampling version of Q, including who
observes and pays each cost, when an agent may exit, and the payoff after a finite run. That model
would then need a new finite-budget incentive or accuracy result, or a stronger impossibility
result, before the pair supports a further claim.

\bibliographystyle{plainurl}
\bibliography{references}

\end{document}
